% TIII_q2, note below the table. Generated by analysis/paper_assets_isc.py; do not edit.
A structure is one target with one seed and five samples. Each time is the median over the runs that compiled the model from scratch, for inputs padded to 256 or 1024 tokens. A process also loads the model and writes its outputs, and its inference time already includes the compilation. Cost is the process wall time per structure times the on-demand list price, unless marked Spot. The amortized cost spreads one compilation over the 8 targets of bucket 1024, as the plan defines it. The warm cost comes from a new process that loads the model compiled by the first run from a saved cache, measured on seed~1 of the 10 pilot targets (2 of them at 1024). For the L4 Triton and the H100, a separate measurement ran the same warm reruns on 10 pilot targets. We read list prices on 8 and 9 October 2026 (A100 on 16 October, H100 on 27 October). They exclude VM setup, idle time, and boot disk. Every main-study run was on demand. Memory is the compiled peak plus the arguments on the device, as the plan defines it (M7). On GPUs, XLA's compiled peak already includes the arguments, so this metric counts them twice there. We found the memory limit by padding the input 7D5C to ever larger buckets. The column gives the largest bucket that ran and the first that failed. The A100 and H100 limits come from a separate measurement. The A100 and H100 rows are follow-ups. The CPU ran 1 target at 256 tokens and 5 at 512, with seeds 1 and 2. The inference time of a fresh process includes the compilation. The amortized cost removes 7/8 of each process's compile time from its wall time before the median is taken, so it can fall below the cost of the median inference time.
